\section{background-covariant tensor current and the Hilbert stress tensor}\label{app:matter-diffeomorphism-current}

consider a local action for tensor fields on an arbitrary background metric

\begin{align}
S_{\mathrm{m}} & =\int _{M}\mathrm{d}^{3}x\sqrt{-g}\,\mathcal{L}_{\mathrm{m}}.
\end{align}

varying the matter action with respect to the metric $\displaystyle{g_{\mu \nu}}$ and the matter field $\displaystyle{\Phi ^{A}}$ gives

\begin{align}
\delta S_{\mathrm{m}} & =\int _{M}\mathrm{d}^{3}x\sqrt{-g}\left(\dfrac{1}{2}T^{\mu \nu}\delta g_{\mu \nu}+E_{A}\delta \Phi ^{A}-\nabla _{\mu}\theta _{\mathrm{m}}^{\mu}[\delta \Phi,\delta g]\right).
\end{align}

using the diffeomorphism covariance and variational setup of Sec.~\ref{sec:theory-setup}, we write the diffeomorphism transformation generated by an arbitrary vector as

\begin{align}
X_{\xi} & =\int \mathrm{d}^{3}x \left( \mathcal{L}_{\xi}g_{\mu \nu} \dfrac{\delta}{\delta g_{\mu \nu}}+ \mathcal{L}_{\xi}\Phi ^{A} \dfrac{\delta}{\delta \Phi ^{A}} \right)
\end{align}

its action on the metric $\displaystyle{g_{\mu \nu}}$ and the matter field $\displaystyle{\Phi ^{A}}$ is

\begin{align}
X_{\xi}\cdot \delta g_{\mu \nu} & =\mathcal{L}_{\xi}g_{\mu \nu}=\nabla _{\mu}\xi _{\nu}+\nabla _{\nu}\xi _{\mu} \\
X_{\xi}\cdot \delta \Phi ^{A} & =\xi ^{\mu}\nabla _{\mu}\Phi ^{A}+(\nabla _{\mu}\xi ^{\nu})(\Delta ^{\mu}_{~\nu}\Phi)^{A}
\end{align}

where $\displaystyle{\Delta ^{\mu}_{~\nu}}$ is the tensor-representation generator appropriate to $\displaystyle{\Phi ^{A}}$. acting $\displaystyle{X_{\xi}}$ on the action gives

\begin{align}
X_{\xi}\cdot \delta S_{\mathrm{m}} & =\int _{M}\mathrm{d}^{3}x\sqrt{ -g }\left(\dfrac{1}{2}T^{\mu \nu}X_{\xi}\cdot \delta g_{\mu \nu}+E_{A}X_{\xi}\cdot \delta \Phi ^{A}\right) \\
 & +\int _{\Sigma _{f}-\Sigma _{i}}\mathrm{d}^{2}x\sqrt{ \sigma }\tau _{\mu}X_{\xi}\cdot \theta ^{\mu}_{\mathrm{m}}[\delta \Phi,\delta g] \\
 & =\int _{M}\mathrm{d}^{3}x\sqrt{ -g }\left(T^{\mu \nu}\nabla _{\mu}\xi _{\nu}+E_{A}(\xi ^{\mu}\nabla _{\mu}\Phi ^{A}+\nabla _{\mu}\xi ^{\nu}(\Delta ^{\mu}_{~\nu}\Phi)^{A})\right) \\
 & +\int _{\Sigma _{f}-\Sigma _{i}}\mathrm{d}^{2}x\sqrt{ \sigma }\tau _{\mu}X_{\xi}\cdot \theta ^{\mu}_{\mathrm{m}} \\
 & =\int _{M}\mathrm{d}^{3}x\sqrt{ -g }\xi ^{\nu}\left(-\nabla _{\mu}T^{\mu}_{~\nu}+E_{A}\nabla _{\nu}\Phi ^{A}-\nabla _{\mu}(E_{A}(\Delta ^{\mu}_{~\nu}\Phi)^{A})\right) \\
 & +\int _{\Sigma _{f}-\Sigma _{i}}\mathrm{d}^{2}x\sqrt{ \sigma }\tau _{\mu}(X_{\xi}\cdot \theta _{\mathrm{m}}^{\mu}-\xi _{\nu}T^{\mu \nu}-E_{A}\xi ^{\nu}(\Delta ^{\mu}_{~\nu}\Phi)^{A})
\end{align}

diffeomorphism covariance of $\displaystyle{S_{\mathrm{m}}}$ also gives

\begin{align}
X_{\xi}\cdot \delta S_{\mathrm{m}} & =\alpha _{\xi}|_{\Sigma _{f}}-\alpha _{\xi}|_{\Sigma _{i}} \\
\alpha _{\xi} & =-\int _{\Sigma}\mathrm{d}^{2}x\sqrt{ \sigma }\tau _{\mu}\xi ^{\mu}\mathcal{L}_{\mathrm{m}}
\end{align}

comparing the two expressions for $\displaystyle{X_{\xi}\cdot \delta S_{\mathrm{m}}}$ gives the bulk and boundary identities

\begin{align}
\nabla _{\mu}T^{\mu}_{~\nu} & =E_{A}\nabla _{\nu}\Phi ^{A}-\nabla _{\mu}(E_{A}(\Delta ^{\mu}_{~\nu}\Phi)^{A}) \\
\nabla _{\mu}\left(X_{\xi}\cdot \theta ^{\mu}_{\mathrm{m}}+\xi ^{\mu}\mathcal{L}_{\mathrm{m}}\right) & =\nabla _{\mu}\left(\xi _{\nu}T^{\mu \nu}+E_{A}\xi ^{\nu}(\Delta ^{\mu}_{~\nu}\Phi)^{A}\right)
\end{align}

imposing the equations of motion $\displaystyle{E_{A}\approx 0}$ gives the on-shell conservation of the Hilbert stress tensor

\begin{align}
\nabla _{\mu}T^{\mu \nu} & \approx 0
\end{align}

locally, the boundary identity is represented by an antisymmetric superpotential $\displaystyle{U_{\xi}^{\mu \nu}}$ as

\begin{align}
X_{\xi}\cdot \theta ^{\mu}_{\mathrm{m}}+\xi ^{\mu}\mathcal{L}_{\mathrm{m}} & =\xi _{\nu}T^{\mu \nu}+E_{A}\xi ^{\nu}(\Delta ^{\mu}_{~\nu}\Phi)^{A}+\nabla _{\nu}U_{\xi}^{\mu \nu}
\end{align}
